% !TeX root = ../papers/nc-p3.tex
\section{二切断の交差階数が同時に一になる繊維の分類}
\label{sec:klein-double-exchange-fourier}
\begin{theorem}\label{thm:full-thin-double-exchange}
節\ref{sec:klein-elliptic-completion}の全薄係数非零の表示で $z=1$, $\lambda=1$ とする。
全ての標準Koszul左右AS正則窓は、位置ごとの文字置換により次のいずれかである。
\begin{enumerate}
\item 定理\ref{thm:klein-singular-exchange-corners}の正表示（最終文字を1だけ移して総文字を零にしたもの）。
\item $g_0,g_1,r_0,r_1\ne0$, $U=r_0+r_1,V=g_0+g_1$, $k=1-UV\ne0$ を取る。
二文字ずつ両半分に属する48語の係数は1とし、全て低位の八語の係数は $U$、
全て高位の八語の係数は $V$ とする。ただし $0000,1111$ の係数にはそれぞれ
$k/g_0,k/g_1$ を、$2323,3232$ にはそれぞれ $k/r_0,k/r_1$ を加える。
\end{enumerate}
記載の全点が正則である。$\lambda=-1$ の全解はこれらの非自明コサイクル捻りで得られる。
四つの族は全点で最小周期1である。
\end{theorem}
\begin{proof}[二位置ずつのFourier変換による被覆]
総文字零を課し、同じ半分の四文字の下位ビットの符号を $e_i\in\{1,-1\}$ とする。
$\prod e_i=1$ なので、低位八係数は一意に
\[
 m+\sum l_i e_i+A e_0e_1+B e_1e_2+C e_0e_2
\]
と書ける。高位は $n,h_i,A',B',C'$ と書く。切断 $01|23$ で $e_0e_1=e_2e_3=s$ とする。
低位角の二行・二列を定数と符号の基底で読んだ係数行列は
\[
 \begin{pmatrix}m+sA&l_2+s l_3\\l_0+s l_1&C+sB\end{pmatrix}.
\]
高位角も同じで、交差ブロックは左上の成分1のみである。切断 $12|30$ では
位置を巡回して $A,B$ を交換する。いずれかの角の右下成分が非零なら、
その成分をピボットに取った残り三行の階数は1以下である。
二つの交差成分が1なので、相手の角は左上の一成分だけとなり、それは非零である。
特に $(C+sB)(C'+sB')=(C+sA)(C'+sA')=0$ が両符号で成り立つ。

一方の半分の全八係数が定数なら正則性を失う。その半分の二文字の和で左右いずれから
評価しても、他方の変数の各半分で係数が一定になる。各奇位置でこの和の一次元を、
各偶位置で二つの半分の差の張る二次元を双対矢空間に取れば、全ての制限関係が零である。
交互にこれらの空間を双対とする経路代数への商で指数増大する。

$C,C'$ がともに非零なら $A/C,B/C=\pm1$、相手の比は逆の符号である。
全位置の下位ビットの平行移動（四移動量の和は零）で $A=B=C$, $A'=B'=-C'$ とする。
低位にピボットを持つ二切断の条件から $(h_0,h_1,h_2,h_3)=(y,-y,y,-y)$、
高位のピボットから $l_i=x$ を得る。もし $C\ne0,C'=0$ なら $A'=B'=0$。
どちらかの切断で低位の両符号にピボットがあれば $h_i=0$ が全てで成立し、高位が定数となる。
従って再び $A=B=C$ にでき、$h=(y,-y,y,-y)$ で $y\ne0$ である。
低位のピボットのない符号では高位角の二非対角成分がともに非零で右下が零だから可逆である。
Schur補完により低位の角は定数であり、やはり $l_i=x$ である。半分を交換すれば逆の場合も従う。
得られた型では低位係数は二つの定数語を除き一定、高位係数は二つの交代語を除き一定となる。

$C=C'=0$ で $A,B,A',B'$ のうち非零があるとする。低位の対係数が一つでも非零なら、
対応する切断で両符号にピボットがあり、高位の全一次係数が零になる。高位の対係数も
非零でなければ定数である。非零なものは $AA'=BB'=0$ により切断の異なる組で、
この場合低位の一次係数も全て零となる。例えば $A,B'\ne0$ なら $B=A'=0$、
二切断条件から $m(n\pm B')=1$ が両方に必要で矛盾する。

残る全対係数零の場合、切断 $ij|kl$ の符号 $s$ に対して
$u=l_i+s l_j,v=l_k+s l_l$, $p=h_i+s h_j,q=h_k+s h_l$ と置く。
三次小行列の条件は $uq=pv=nuv=mpq=0$ で尽きる。
どちらの一次係数も全ては零でない。隣接する $l_i^2,l_j^2$ が異なれば対向する組の
$h_k,h_l$ はともに零である。$h$ の支持が隣接二位置ちょうどなら、その両端と外側との
二つの隣接対で平方が異なるので逆に $l$ 全体が零になる。
一位置だけなら $l$ の支持もその位置だけである。この場合、その位置を含まない連続二辺では
各厚色の係数空間が、各半分で一定なベクトルの二平面になる。
両方の関係空間に $(e_0+e_1)\otimes(e_0+e_1)$ が入り、
定理\ref{thm:no-consecutive-pure-relations}で排除される（ここで $e_i$ は文字基底）。
従って全 $l_i^2$ は同一非零であり、$h_i^2$ も同様である。
平行移動後 $l=(x,x,x,\delta x)$, $\delta=\pm1$ とすると
$h=(y,-\delta y,y,-y)$。$\delta=-1$ は
$(X,Y,W,V)=(m+2x,m-2x,n+2y,n-2y)$ の第一型である。
$\delta=1$ は先の定数語・交代語だけが異なる型に含まれる（$nuv=mpq=0$ から $m=n=0$）。

第二型の四例外語への加算量を順に $p,q,j,t$ と置くと厚い階数の上界は
$Vpq=(1-UV)(p+q),\ Ujt=(1-UV)(j+t)$ である。
これは定数値 $U,V$ を持つ二行対に対する三次小行列式そのものである。
$k\ne0$ なら片方の加算量だけ零にはならず、二つとも零なら定数半分として除かれる。
従って $g_0=k/p,g_1=k/q,r_0=k/j,r_1=k/t$ が記載の表示を与える。
$k=0$ なら $pq=jt=0$ で、両方の低位加算量が零なら第一厚色の階数が1に落ちる。
一方だけ非零ならその二列の差が文字の平方の純関係となる。
テンソルは偶位置の交換にも奇位置の交換にも対称であるため、
補題\ref{lem:bisymmetric-reducible-quick-test}の評価判定の必要条件、または連続純関係の排除で落ちる。
\end{proof}

\begin{proof}[十分性と周期]
第二型の色0関係の係数平面は $(a,b,g_0a+g_1b,g_0a+g_1b)$、
色1は $(r_0c+r_1d,r_0c+r_1d,c,d)$ である。薄い二色は全成分1。
各々を $ij$ の係数行列にまとめると交差ブロックはともに
$\left(\begin{smallmatrix}u&v\\v&u\end{smallmatrix}\right)$。
非零の交差を持つ純関係なら $v=\pm u$ であり、同じ符号の二文字ベクトルが左右各半分に現れる。
従って $a=b$, $c=d=\pm Va$, $a=UVa$ となり矛盾する。交差が零なら片方の角が零で、
四係数 $g_0,g_1,r_0,r_1$ の非消滅により残る角も階数1にはならない。

評価点を $(x_0,x_1,x_2,x_3)$ とし $S=x_2y_2+x_3y_3,T=x_0y_1+x_1y_0$ と置く。
厚い四式は $x_0y_0=-g_0S,x_1y_1=-g_1S,x_2y_3=-r_0T,x_3y_2=-r_1T$。
全座標が非零なら $S$ が $y$ を決める。低位の一方だけが零なら $S=0$ となり、
高位二座標が非零のとき残りの自由度は高々一つである。高位にも一つ零があれば $T=0$ で、
四式と $S=T=0$ が核を零にする。高位だけ一座標が零の場合も $T$ から先に同じ議論をする。
片半分だけが非零なら厚い式が同じ半分の二座標を消し、薄い式が少なくとも一条件を加える。
右評価は $r_0,r_1$ の交換で同じである。反対切断の薄い二色の各核は
二つの交差積の差を含み、その和が二つの一次式の積になる。よって
補題\ref{lem:bisymmetric-reducible-quick-test}により全次数で正則である。

$P$ を高位二文字の交換とすると関係六平面は $R^{\mathsf t}=(P\otimes P)R$ を満たす。
全辺が $R,R^{\mathsf t}$ を反復するので最小周期は1である。
非自明捻りで $P$ は次数の自己同型 $i\mapsto i+(i\mathbin{\mathrm{div}}2)$ を誘導し、
二コサイクルの比は $(-1)^{i_{\rm hi}j_{\rm hi}}$、すなわち
$\beta(i)=(\sqrt{-1})^{i_{\rm hi}}$ の余境界である。よって一文字倍率で補正した $P$ が
捻り後も一位置移動を実現する。さらに全薄い48語上で
$\prod_{i<j}\mu(q_i,q_j)$ に対角文字倍率を掛けると、ちょうど $\lambda$ を $-\lambda$ にする。
低位の厚角 $A,B$ への作用は $DAD,-DBD$ ($D=\operatorname{diag}(1,-1)$)。
\path{computations/verify_double_exchange_fourier.py} が全語でこの符号と切断の式を確かめる。
全次数の正則性は捻りで保存される。
\end{proof}

\begin{theorem}\label{thm:double-exchange-frames}
第二型は正表示で反対階数8、捻りで
\[
 16-[U=0]-[V=0]-[k^2+(4Ug_0g_1+kV)(4Vr_0r_1+kU)=0]
\]
を持つ。両表示とも印を忘れた添字保持同型の必要十分条件を、全パラメータで次の有限な枠比較により得る。
正表示では各反対核にある四つの純テンソルの因子を枠とする。捻りでは各位置の被約二十点から
枠を取り、座標平面を含む全ての候補を調べればよい。四位置の枠を順序付け、テンソルの
非零台とその整数文字関係の単項式比を記録する。正規形としてこの有限集合を取れば、
二集合が交わることと同型であることは同値である。第一型と第二型および各々の正負の間に同型はない。
第二型の正負の全点が前節の混合補完族の特殊化となる。
\end{theorem}
\begin{proof}[反対核と捻りの二十点]
正表示の反対切断は色順に階数 $(4,2,1,1)$ である。色0の行列式は
$-k^3/(g_0g_1r_0r_1)$、色1は二組で行・列が等しく縮約行列の行列式が $UV-1$ となる。
対称二次式の核は $e_0 e_2-e_1 e_3,e_0e_3-e_1e_2$ の鉛筆で、
積となる二元は $(e_0\pm e_1)(e_2\mp e_3)$ である（符号を入れ替えても同じ二組）。
二つの反対位置を別空間としても、純テンソルの対称化はその二因子の積であり、
従って全反対核の純テンソルはこれらの各順序の四点で尽きる。
左右の四因子は独立で、相手側の反対切断にも同じ議論を適用できる。
捻りの薄い反対二色は行列式 $-16$、色1は二組の和と差に分けると階数
$2+[U\ne0]+[V\ne0]$ である。色0の二つの角は可逆、交差は階数1であり、
Schur補完の階数損失は高々1である。行列式を計算すれば上記の式を得る。

捻りの左点を斉次座標 $(a,b,c,d)$ で書く。二次関係から
$S=c y_2-d y_3$, $T=a y_1+b y_0$ とおいて
\[
 ay_0=-g_0 S,\quad by_1=-g_1 S,\quad cy_3=r_0T,\quad dy_2=-r_1T .
\]
薄い二式は $ay_2+by_3+cy_0-dy_1=ay_3+by_2-cy_1+dy_0=0$。
全座標が非零なら $z=b/a,w=d/c$, $G=g_1/z+g_0z$, $H=r_1/w+r_0w$ に対し
$T=-GS,S=-HT$ で $GH=1$。薄い二式は $G,c^2/a^2$ の同次二式となるが係数行列式は
\[
 (r_1/w-r_0z)(g_1/z-g_0w)-(-r_0+r_1z/w)(-g_0+g_1w/z)
 =GH-UV=k\ne0
\]
となり矛盾する。一座標が零で他の三つが非零の場合は次の四組を得る。
\[
\begin{array}{c|cc}
d=0 & a^2=g_0c^2& b^2=-g_1c^2\\
c=0 & a^2=-g_0d^2& b^2=g_1d^2\\
b=0 & c^2=r_0a^2& d^2=-r_1a^2\\
a=0 & c^2=-r_0b^2& d^2=r_1b^2 .
\end{array}
\]
二座標以上が消える場合は $(1,\pm1,0,0),(0,0,1,\pm1)$ のみである。
実際片半分が零なら厚い式が同じ半分の $y$ を零とし、残りの二式の行列式は
二つの平方の差となる。両半分に一零ずつあれば $S=T=0$ から四座標とも消える。
これで相異なる二十点があり、有限な期待余次元の最大小行列式スキームの長さ20から被約である。
他の位置も既に構成した一位置同型で同じとなる。

各座標平面はその上の四つの片零点で張られ、そこに純半分の二点もある。
各位置で配置の独立三点が張り六点以上を含む平面を全て取り、その中の法線が
独立な四平面を選べば、座標四面体を含む固有な有限集合が得られる。
追加の六点平面を持つ特殊値もこの規則に含む。どちらの符号でも残る対角倍率は
非零語の文字指数行列の飽和整数核で正確に比較できる。
正表示二型は反対階数6と8で異なり、負表示はいずれもそれより大きい。
負表示二型は点配置の曲線部分の有無で異なる。
\end{proof}

\begin{proof}[混合補完からの特殊化]
第二型の低位角を $A=\left(\begin{smallmatrix}a&b\\b&d\end{smallmatrix}\right)$ として
$b=U,a=b+k/g_0,d=b+k/g_1$ と取る。$\det A=k/(g_0g_1)$ である。
$\lambda=1+\epsilon,j=r_0-r_1,f=b+\epsilon j$ とし、前節の式で $t,g$ を決める。
$tb-1$ の剰余は $-k$、$g$ の剰余は $b$ となり、$\det B$ は正確に位数2を持つ。
実際その主要係数は $(b-j)(b+j)\det A/((a-b)(b-d))$ である。
$T'$ の非対角の剰余は $V$、対角の剰余は順に $V+2k/(b+j),V+2k/(b-j)$ となる。
従って第二型の四例外語が全て希望の値を持つ。一般繊維では
$x=ad-fg$ の剰余も $\det A\ne0$ であり、前節の混合補完条件に属する。
コサイクルの符号補正はこの補完で $b$ と $\lambda$ の符号を変え、他の $a,d,f,g$ を保つ。
よって負表示も $\lambda=-1$ での同じ族の特殊化である。
\end{proof}
